\chap{Abstract} 


%%%% your abstract goes here (word limit: 350)
% \Blindtext[3]
Groups of cells, including clusters of cancer cells, developing organs, and whole organisms, can both grow and break apart.
What controls when cells detach from a group, and what sets the size of the pieces that break off?
How accurately can a cell sense that it has touched a neighbor?
This dissertation develops simple theoretical and computational models for two problems: the mechanics of cell-cell detachment and tissue fracture, and the physical limits on cell-cell sensing.

I first study how cancer cells dissociate from a collectively invading strand in confinement, an \textit{in vitro} mimic of cancer metastasis.
Using a phase-field model, I find that single cells break off most often, that fast ``leader'' cells explain this, and that stronger cell-cell adhesion produces fewer but larger ruptures.
Unjamming the tissue is necessary but not sufficient for rupture.
Building on this, I ask what sets the size of a growing tissue, where division competes with motility-driven fracture.
Modeling it as self-propelled particles joined by breakable junctions, I find that its size depends only on the ratio of the junction break rate to the cell division rate, so motility can effectively regulate organ size.

To follow rupture into larger tissues, such as \textit{Trichoplax adhaerens}, I need a model that scales to many cells while keeping tissue fluidity.
I resort to the finite Voronoi model, a variant of the standard Voronoi model for nonconfluent tissue, and find that, surprisingly, its detachment forces diverge as cells lose contact, so the fracture timescale unphysically depends on the simulation time step.
I then introduce a regularization, calibrated against a deformable polygon model.

Finally, I turn from mechanics to sensing.
In contact inhibition of locomotion, a cell repolarizes away from a neighbor upon contact.
Modeling this as a cell reading ephrin bound to Eph receptors against nonspecific interfering ligands, I find that a cell can reliably decide whether a neighbor is present, even when the interfering ligand concentration is ten times that of the correct ligands.
Together, these models show how cell groups hold together, break apart, and sense their neighbors.


%% list of keywords seperated by comma
%\keywords{Johns Hopkins, PhD, Masters, dissertation, thesis \LaTeX, template.}


%%%%  committee members (add it right after the abstract w/o page break)
\begin{singlespace}

    %% if you have co-advisor, add here w/ \vspace{0.1in} as shown below
    %% alternatively you can use minipage environment to put side-by-side
    \section*{Primary reader and thesis advisor}
    
    Dr. Brian A. Camley \\
    %Professor\\
    Department of Physics and Astronomy\\
    Johns Hopkins University, Baltimore, MD 


    \section*{Secondary readers}
    
    Dr. Stewart Hawking\\
    Professor\\
    Department of Biology \\
    Johns Hopkins University, Baltimore, MD 
    
    \vspace{0.1in}
    
    Dr. Jimmy Watson \\
    Professor\\
    Department of Biology \\
    Johns Hopkins University, Baltimore, MD 

    %% you can add more readers if you have them on your committee 
    %% use \vspace{0.1in} in between members for clarity
    %% you can also place committee members side-by-side using `minipage`

    \section*{Thesis Committee}

    \begin{minipage}[t]{0.45\textwidth}
    Dr. Brian A. Camley \\
    Department of Physics and Astronomy\\
    Johns Hopkins University, Baltimore, MD 
    \end{minipage}
    \hfill
    \begin{minipage}[t]{0.45\textwidth}
    Dr. Daniel H. Reich \\
    Department of Physics and Astronomy\\
    Johns Hopkins University, Baltimore, MD 
    \end{minipage}

    \vspace{0.1in}

    \begin{minipage}[t]{0.45\textwidth}
    Dr. Yaojun Zhang \\
    Department of Physics and Astronomy\\
    Johns Hopkins University, Baltimore, MD 
    \end{minipage}


\end{singlespace}